% =============================================================================
% Chapter 6  Polynomial Nonlinearity  --  ~4 pp
% rubric: Project Body (Modelling) -- Model Formulation /30.
% Plan: ThesisPlanning.md §5.6 (supporting contribution S4).
%
% PURPOSE     Answer Q(iii): separate what the doubled space can claim (an
%             exact encoding, a projection) from what the Carleman lift can (a
%             preparation, in both forms), and set up the fluids.
% SOURCES     ThesisPlanning.md §3.6, §5.3 (R6-R10), §6.5, §7.7;
%             pihm/docs/02_methods.md §3.4-3.5; claims A-10, B-12, B-18, B-19,
%             B-22, C-09, C-10; pihm/src/pihm/carleman.py, problems/fluids.py,
%             circuits/tensor.py, circuits/folds.py.
% WRITE WHEN  Gate G5 (wave W4); §6.3 at gate G6 (wave W5).
% =============================================================================

\subsection{Two Treatments}
\label{sec:nonlinear-overview}
% ~0.25 pp
% LAND: nonlinearity is the second axis of sec:pipeline-inputs; the doubled
%   space (the paper's; run here in the standard form only, since it projects
%   in either) and the Carleman lift (both forms). One sentence tying back to
%   sec:descriptor-idea's structural principle (the same idea applied again,
%   not new machinery).
% REVISE: the LAND used to say "both treatments work in both forms" -- the
%   doubled space is run in the standard form only (CLAUDE.md Standing Facts,
%   Nonlinearity): it projects in either form, so the thesis never builds it in
%   the descriptor form. Only the Carleman lift runs in both.

\subsection{The Doubled Space: Exact Encoding, Degenerate Kernel}
\label{sec:doubled-space}
% ~0.6 pp
% LAND, in order:
%   1. Wu et al.'s construction (sec:pihm-doubled), run in the standard form
%      ONLY: an EXACT encoding (claim B-12 for the fold).
%   2. The kernel dimension: >= N^2 - 2N measured (the paper states
%      >= 2^{2n-1}; claim A-10).
%   3. prop:no-penalty: QITE returns the projection of phi_0 onto the kernel;
%      no Hamiltonian enforces the product structure (a kernel containing
%      psi (x) psi and phi (x) phi contains their sum, which is not a product
%      state), so the kernel is one-dimensional only if the construction singles
%      out one solution in advance, which an answer-independent penalty cannot
%      do. Proof -> Appendix H (app:proof-no-penalty).
%   4. The claim allowed: "the encoding is exact" and "QITE projects onto the
%      kernel". NOT: "prepares the solution".

The doubled Hilbert space method can create degenerate kernels. \Gls{qite} can't select the physical solution from amongst a set of degenerate ground states, so although it is possible to encode the solution to the original problem in the ground state of the Hamiltonian, it is not possible to extract it from the ground state of the descriptor Hamiltonian.
% REVISE: moved here from the draft's descriptor chapter. The last clause
%   should read "doubled-space Hamiltonian", not "descriptor Hamiltonian"; and
%   "can't" breaks the register (Style Trap 5).

\begin{proposition}[Imaginary time projects onto a degenerate kernel]
\label{prop:no-penalty}
% STATE: for the positive semi-definite doubled-space Hamiltonian H_(x) and any
%   initial state with Pi_ker phi_0 != 0, the normalised imaginary-time state
%   converges to the normalised projection of phi_0 onto the kernel. SCOPE of the
%   corollary: the kernel is a linear subspace, so one containing psi (x) psi and
%   phi (x) phi contains their sum, which is not a product state; no Hamiltonian
%   therefore enforces the product structure, and an answer-independent penalty
%   cannot make the kernel one-dimensional. (A penalty that already knows the
%   solution can, and is excluded: it is answer-dependent.)
\begin{equation}
    \lim_{\beta\to\infty}\frac{e^{-\beta H_{\otimes}}\,\varphi_0}{\bigl\|e^{-\beta H_{\otimes}}\,\varphi_0\bigr\|}
    = \frac{\Pi_{\ker H_{\otimes}}\,\varphi_0}{\bigl\|\Pi_{\ker H_{\otimes}}\,\varphi_0\bigr\|},
    \qquad
    \Psi = \psi \otimes \psi .
    \label{eq:doubled-projection}
\end{equation}
\end{proposition}

\subsection{The Carleman Lift in Space--Time}
\label{sec:carleman}
% ~1.25 pp
% LAND, in order:
%   1. The lift (eq:carleman-lift) and its truncation error O(rho^K), with
%      rho < 1 STATED AS A SUFFICIENT CONDITION, A WORST CASE OVER ALL STATES
%      (eq:carleman-ratio) -- NOT necessary; the convergence is measured per
%      order against the exact lift (sec:results-nonlinear; on Burgers it
%      follows a/nu, not rho).
%   2. A linear IVP has a unique solution, so the space-time residual has a
%      ONE-DIMENSIONAL near-kernel -- a real preparation claim (claim C-09).
%   3. The two time discretisations (eq:carleman-spacetime): the same kernel,
%      differing only by a banded invertible left factor -- why R8-R10 compare
%      like for like. The descriptor's time axis lowers the degree.
%   4. Initial data enter through ratio rows, F (x) <tau(-1)| (tab:constraint-kinds).
%   5. Uniqueness is Carleman's, not the descriptor's (claim B-22 retired).
% OFFLOAD: the symmetric and tensor layouts, the ratio rows and the space-time
%   forms -> Appendix H (app:nonlinear).

\begin{equation}
    \frac{du}{dt} = \sum_{p=1}^{P} F_p\,u^{\otimes p}
    \quad\longrightarrow\quad
    \frac{dz}{dt} = \mathbb{A}_K\,z,
    \qquad
    z = \bigl(u,\ u^{\otimes 2},\ \dots,\ u^{\otimes K}\bigr),
    \label{eq:carleman-lift}
\end{equation}
\begin{equation}
    \rho = \frac{\|u_0\|\,\|F_2\|}{\bigl|\operatorname{Re}\lambda_1(F_1)\bigr|} < 1,
    \qquad
    \text{truncation error} = O\bigl(\rho^{K}\bigr),
    \label{eq:carleman-ratio}
\end{equation}
\begin{equation}
    \begin{aligned}
        A_{\mathrm{st}}^{\mathrm{std}} &= \G_t \otimes I - \tfrac{T}{2}\, I \otimes \mathbb{A}_K, \\
        A_{\mathrm{st}}^{\mathrm{desc}} &= \Dt_t \otimes I - \tfrac{T}{2}\, \Bt_t \otimes \mathbb{A}_K
        = \bigl(\Bt_t \otimes I\bigr)\, A_{\mathrm{st}}^{\mathrm{std}},
    \end{aligned}
    \qquad
    \xi = \frac{2t}{T} - 1 .
    \label{eq:carleman-spacetime}
\end{equation}

\subsection{Products and Couplings on a Circuit}
\label{sec:folds}
% ~0.75 pp
% LAND, in order:
%   1. prop:fold-optimal: in the nodal basis a product is pointwise; 0 ancilla,
%      (p - 1)(n + l) CNOTs, and alpha = 2 = ||n_1|| for p = 2 -- OPTIMAL, not
%      merely tight (operator-verified to ~1e-15). Define l (the output
%      register's growth) where the fold is introduced.
%   2. The Fourier fold: QFT, CNOT copy, QFT^dagger -- native and exact.
%   3. The Chebyshev DCT is budgeted as a cited primitive (klappenecker2001dct;
%      decision 6); NOT built gate-level, said plainly.
%   4. The lift's coupling A_K: as matchings in the symmetric layout, or in the
%      TENSOR LAYOUT WITH 2K - 1 TERMS (not K^2), each F_p on adjacent slots,
%      which reaches field scale.
%      MOVED here from 5_descriptor.tex (2026-10-06): a large irregular coupling is not expanded in Pauli
%      words but split into matchings, each with at most one entry in every row and column; a greedy
%      colouring uses at most 2*Delta - 1 of them, Delta the largest row degree.
% OFFLOAD: proof -> Appendix H (app:proof-fold-optimal); the symmetric and
%   tensor layouts in full -> Appendix H (app:nonlinear).

\begin{proposition}[The nodal fold is optimal]
\label{prop:fold-optimal}
% STATE: the degree-p product fold factorises through the Chebyshev--Gauss
%   transform as below, needs no ancilla, and for p = 2 its subnormalisation
%   equals its operator norm.
\begin{equation}
    \nfold_1^{(p)} = 2^{p\ell/2}\,\Ucg^{\dagger}\,\Delta_p\,\Ucg^{\otimes p},
    \qquad
    \alpha\bigl(\nfold_1^{(2)}\bigr) = 2 = \bigl\|\nfold_1^{(2)}\bigr\| .
    \label{eq:nodal-fold}
\end{equation}
\end{proposition}

\subsection{Fourier--Galerkin Fluids: Burgers and Navier--Stokes}
\label{sec:fluids}
% ~0.75 pp
% LAND, in order:
%   1. Burgers (R9, eq:burgers): Fourier-Galerkin with |k| <= K_c,
%      F_1 = diag(-nu k^2), F_2 the (u^2)_x / 2 convolution; (a, nu) chosen so
%      the measured rho <= 0.5.
%   2. 2-D vorticity Navier-Stokes (R10, eq:navier-stokes), psi_hat =
%      omega_hat / |k|^2 eliminated; exact rational couplings.
%   3. Taylor-Green (its modes share |k|, so the lift is exact: a validation)
%      and a counter-rotating vortex pair (active triads).
%   4. References: the Galerkin system's IVP, Cole-Hopf for Burgers, a
%      pseudo-spectral simulation for the truncation's own error.
% OFFLOAD: F_1 and F_2 for Burgers and NS in full -> Appendix H.
% REVISE: this LAND used to put the tensor-layout coupling here as "K^2
%   structured terms"; it is 2K - 1 terms (not K^2), and the construction
%   itself has moved to sec:folds point 4, since it is a circuit-level coupling
%   the fold section now covers, not a fluids-specific fact.

\begin{equation}
    u_t + u\,u_x = \nu\,u_{xx},
    \qquad
    F_1 = \operatorname{diag}\bigl(-\nu k^2\bigr),
    \label{eq:burgers}
\end{equation}
\begin{equation}
    \omega_t + J(\psi, \omega) = \nu\,\nabla^2\omega,
    \qquad
    \nabla^2\psi = -\omega \ \ \text{on } [0, 2\pi)^2,
    \qquad
    \omega_0^{\mathrm{TG}} = 2\sin x\,\sin y .
    \label{eq:navier-stokes}
\end{equation}

\subsection{Preparation Versus Projection}
\label{sec:prep-vs-proj}
% ~0.4 pp
% LAND: TABLE T8: which route sits on which side, and why. THE CAVEAT THAT
%   MATTERS: a unique, trusted kernel certifies the lifted LINEAR problem, not
%   the lift's convergence, which is measured separately
%   (sec:results-nonlinear). One sentence tying back to sec:descriptor-idea:
%   the mass form is the same banded left factor, applied in time.
% TAB: tab:prep-vs-proj (T8).

\begin{table}[htbp]
    \centering
    \caption{Which nonlinear route prepares the solution, and which only projects onto a kernel that contains it.}
    \label{tab:prep-vs-proj}
    \small
    \begin{tabular}{@{}p{0.24\linewidth} p{0.22\linewidth} p{0.26\linewidth} p{0.18\linewidth}@{}}
        \toprule
        Route & Kernel & What imaginary time returns & Precondition \\
        \midrule
        Doubled space (Wu et al.; R6) & degenerate, dimension $\ge 2^{2n-1}$ & the projection of $\varphi_0$ onto the kernel & none \\
        Carleman lift (R8--R10), either form & one-dimensional near-kernel & the truncated solution & $\rho < 1$ \\
        \bottomrule
    \end{tabular}
\end{table}
% REVISE: the "Precondition" column's $\rho < 1$ reads as a necessary
%   condition. CLAUDE.md Standing Facts (Nonlinearity) and ThesisPlanning.md
%   §5.6 state rho < 1 is SUFFICIENT, not necessary, and a worst case over all
%   states; the lift's actual convergence is measured per order against the
%   exact lift (on Burgers it follows a/nu, not rho). Either retitle the
%   column ("Sufficient for the O(rho^K) bound") or add a footnote.

% CHECKLIST (marker)
%   1. Is "prepared" never used for the doubled space?
%   2. Is rho < 1 stated as SUFFICIENT, not necessary?
%   3. Is fold optimality distinguished from tightness?
%   4. Is the DCT named as unbuilt, with the reason?
%   5. Is uniqueness attributed to Carleman?
%   6. Is trust distinguished from convergence (sec:prep-vs-proj)?
% TRAPS  "K^2 terms" (it is 2K - 1); "both forms" for the doubled space (it is
%        standard-form only); rho < 1 as necessary.
